\documentclass[10pt]{article}
\usepackage[margin=2cm]{geometry}
\usepackage{booktabs}
\usepackage{siunitx}
\title{3-D BCH encoder, $N=15$, $T=2$: three architectures after place and route}
\author{}
\date{31 October 2026}
\begin{document}
\maketitle

\section{Designs}
All three encode $K=3163$ message bits into an $N^3=3375$-bit codeword over GF(16).
\begin{description}
\item[Design 1, baseline] (\texttt{bch3d\_enc\_baseline}): Mondal \& Garani 2021, extended to 3-D. Every output point has its own conjugate-class units and XOR array. One codeword per clock.
\item[Design 2, multiplier optimisation] (\texttt{bch3d\_enc\_opt}): compute once, route many. A shared bank of computation-class products (11,500 wires) feeds 3,375 XOR trees. One codeword per clock.
\item[Design 3, time-multiplexed] (\texttt{bch3d\_enc\_tm\_asm}): spectral registers rotated by $\alpha^{-j_3}$ and 225 XOR trees reused over 15 cycles, plus an assembly shift register. One codeword per 15 clocks.
\end{description}
Designs 1 and 2 are wrapped with input and output registers (\texttt{\_reg}, 6,538 FFs, 2-cycle latency). Both wrappers passed RTL simulation (3,170 vectors, 0 errors).

Flow: Vivado 2026.1, out-of-context, \texttt{synth\_design -flatten\_hierarchy rebuilt}, \texttt{opt\_design}, \texttt{place\_design}, \texttt{route\_design}, vectorless power. Machine: 15.3\,GB RAM.

\section{Results}
\begin{table}[h]\centering\small
\begin{tabular}{lrrrr}
\toprule
 & D1 slice (225 pts) & D2 full & D3 Virtex-7 & D3 Artix-7 \\
\midrule
Part & xc7v2000t-1 & xc7v2000t-1 & xc7v2000t-1 & xc7a200t-1 \\
Status & routed & \textbf{not routable} & routed & routed \\
LUTs & 37,882 & 181,728 (synth) & 17,196 & 17,209 \\
FFs & 3,354 & 6,538 & 10,200 & 10,168 \\
Clock period & 10\,ns & 10\,ns & 6\,ns & 8.5\,ns \\
WNS & +1.618 & +0.383 (post-place) & +0.062 & +0.139 \\
Fmax & $\approx$119\,MHz & -- & 166.7\,MHz & 117.6\,MHz \\
Power total / dynamic & 3.03 / 2.36\,W & -- & 3.35\,W & 1.92 / 1.80\,W \\
Peak memory & 5.3\,GB & 9.2\,GB & -- & -- \\
Run time & 11.7\,min & 197\,min & -- & $\approx$10\,min \\
\bottomrule
\end{tabular}
\caption{Measured results. ``--'' means no number exists for that run.}
\end{table}

\paragraph{Design 2.} Synthesis with \texttt{rebuilt} gives 181,728 LUTs (14.9\% of the device), a third of the 544,597 LUTs from \texttt{flatten none}. Placement closes timing at 10\,ns (estimated WNS $+0.383$\,ns), but routing fails: 25,832 signals unrouted, 66,912 node overlaps, peak congestion 99.8\%. Each of the 3,375 XOR trees reads from across the whole 3,163-bit message, so the wiring is all-to-all. The design runs out of routing, not of LUTs.

\paragraph{Design 2 slice.} A standalone 225-point slice (plane $i''=0$, 10,608 bank products, 225 trees) routes at 10\,ns on xc7v2000t-1. It uses 12,816 LUTs and 3,354 FFs, with WNS $+2.332$\,ns ($\approx$130\,MHz) and 1.50\,W (0.85\,W dynamic). The critical path is 7.634\,ns, with 5 LUT levels and 91\% routing. Scaling by 15 gives $\approx$192\,k LUTs, close to the 181,728 LUTs from synthesising the full design. Compared with design 1's slice it uses 3$\times$ fewer LUTs and 2.8$\times$ less dynamic power. In full, design 2 would deliver $\approx$130\,M codewords/s at $\approx$13.4\,W, or $\approx$0.10\,$\mu$J per codeword, if it could be routed.

\paragraph{Design 1.} Design 1 has the same all-to-all wiring and more logic, so it was not placed and routed in full (fallback rule of the plan). The 225-point slice routes at 10\,ns. Scaling it by 15 gives $\approx$568\,k LUTs and $\approx$36\,W dynamic power. These are estimates, and best-case ones: the slice has 15$\times$ less input fan-out than the full design. Since the smaller design 2 already fails to route, the full design 1 is expected to fail as well. A synthesis-only run with \texttt{rebuilt} was stopped after 13.5\,min at 10.2\,GB, during area optimisation, before free RAM ran out. On this 15.3\,GB machine design 1 cannot even be synthesised flattened.

\section{Throughput and efficiency}
\begin{table}[h]\centering\small
\begin{tabular}{lrrrr}
\toprule
 & Codewords/s & Gbit/s & Energy / codeword & Gbit/s per k LUT \\
\midrule
D1 full (estimate) & $\le$119\,M & $\le$402 & $\approx$0.31\,$\mu$J & $\approx$0.71 \\
D2 full (if routed at 10\,ns) & 100\,M & 337.5 & -- & 1.86 \\
D3 Virtex-7, 6\,ns & 11.1\,M & 37.5 & 0.30\,$\mu$J & 2.18 \\
D3 Artix-7, 8.5\,ns & 7.84\,M & 26.5 & 0.25\,$\mu$J & 1.54 \\
\bottomrule
\end{tabular}
\caption{Gbit/s counts codeword bits. Design 1's numbers are slice $\times$15. Design 2's are what it would deliver at the post-place clock, had it routed.}
\end{table}

\section{Conclusions}
\begin{itemize}
\item Design 3 is the only full $N=15$ encoder that routes. It uses about 1/10 of design 2's LUTs and fits a low-cost Artix-7 (12.9\% of the 200T).
\item Design 3 gives up 15$\times$ throughput per clock, but per LUT it is close to or better than designs 1 and 2. Its energy per codeword matches the design 1 estimate.
\item Moving from design 1 to design 2 removes duplicated multipliers, but routing becomes the limit. Time-multiplexing cuts the number of simultaneous XOR trees from 3,375 to 225, and that is what makes the design routable.
\end{itemize}
\end{document}
